Large language models increasingly act as readers, judges, and filters for text, yet most detection work still asks only whether text is human- or machine-written. We study a broader question: can an LLM detect how a text was produced, and does that production mode behave like a latent variable in representation space? We build a controlled six-mode corpus from 60 \hcthree anchors, balanced across open-domain QA, finance, and medicine, with matched variants for human originals, direct LLM generations, LLM-humanized text, human text polished by an LLM, dictated/spoken rewrites, and keyboard-noisy text. We then evaluate zero-shot mode classification with \gptfourone, compare it to a grouped stylometric baseline, probe \embeddingmodel embeddings with grouped cross-validation, and test whether mode-specific embedding deltas are consistent across content.

The results support a narrow but strong claim. On the controlled benchmark, \gptfourone achieves perfect six-way classification, while the stylometric baseline reaches 0.400 accuracy and the grouped embedding probe reaches 0.489 accuracy on unseen content anchors. All five transformed modes also show statistically significant within-mode delta consistency, with the strongest effects for keyboard noise and dictated speech. However, the same prompting setup reaches only 0.256 accuracy on six-way source attribution in \semeval Subtask B and collapses most non-human systems into \chatgptlabel or human. We conclude that current LLMs read broad production modes well, but zero-shot generator identity remains substantially harder.
